\documentclass[12pt,a4paper]{article}
\usepackage[UTF8]{ctex}          % 添加中文支持
\usepackage{amsmath,amssymb,amsfonts,geometry,enumitem}
\geometry{left=2.5cm,right=2.5cm,top=2.5cm,bottom=2.5cm}
\title{2024年新高考Ⅰ卷·数学 参考答案}
\date{}
\begin{document}
\maketitle

\section*{一、单项选择题}

\begin{enumerate}[label=\arabic*.]
\item A
\item C
\item D
\item A
\item B
\item B
\item C
\item B
\end{enumerate}

\section*{二、多项选择题}

\begin{enumerate}[label=\arabic*.]
\item BC
\item ACD
\item BC
\end{enumerate}

\section*{三、填空题}

\begin{enumerate}[label=\arabic*.]
\item $\dfrac{3}{2}$
\item $\ln 2$
\item $\dfrac{7}{16}$
\end{enumerate}

\section*{四、解答题}

\begin{enumerate}[label=\arabic*.]
\item 
\begin{enumerate}[label=(\arabic*)]
\item 由 $a^2+b^2-c^2=\sqrt{2}ab$，根据余弦定理 $\cos C=\dfrac{a^2+b^2-c^2}{2ab}=\dfrac{\sqrt{2}}{2}$，故 $C=\dfrac{\pi}{4}$。

又 $\sin C=\sqrt{2}\cos B$，即 $\dfrac{\sqrt{2}}{2}=\sqrt{2}\cos B$，得 $\cos B=\dfrac{1}{2}$，故 $B=\dfrac{\pi}{3}$。

\item 由 $A=\pi-B-C=\pi-\dfrac{\pi}{3}-\dfrac{\pi}{4}=\dfrac{5\pi}{12}$。

面积 $S=\dfrac{1}{2}ac\sin B=\dfrac{1}{2}ac\cdot\dfrac{\sqrt{3}}{2}=\dfrac{\sqrt{3}}{4}ac=3+\sqrt{3}$，得 $ac=4\sqrt{3}+4$。

由正弦定理 $\dfrac{a}{\sin A}=\dfrac{c}{\sin C}$，得 $a=\dfrac{c\sin A}{\sin C}$。

代入 $ac$ 可解得 $c=2\sqrt{2}$。
\end{enumerate}

\item 
\begin{enumerate}[label=(\arabic*)]
\item 将 $A(0,3)$ 代入椭圆得 $b^2=9$。将 $P(3,\frac{3}{2})$ 代入得 $\dfrac{9}{a^2}+\dfrac{9/4}{9}=1$，即 $\dfrac{9}{a^2}+\dfrac{1}{4}=1$，$\dfrac{9}{a^2}=\dfrac{3}{4}$，$a^2=12$。

$e=\dfrac{c}{a}=\dfrac{\sqrt{a^2-b^2}}{a}=\dfrac{\sqrt{12-9}}{\sqrt{12}}=\dfrac{\sqrt{3}}{2\sqrt{3}}=\dfrac{1}{2}$。

\item 直线 $l$ 的方程为 $y-\frac{3}{2}=k(x-3)$，即 $y=kx-3k+\frac{3}{2}$。

点 $A$ 到直线 $l$ 的距离 $d=\dfrac{|0-3-(-3k+\frac{3}{2})|}{\sqrt{k^2+1}}=\dfrac{|3k-\frac{9}{2}|}{\sqrt{k^2+1}}$。

$|AP|=\sqrt{3^2+(\frac{3}{2}-3)^2}=\sqrt{9+\frac{9}{4}}=\dfrac{3\sqrt{5}}{2}$。

$S_{\triangle ABP}=\dfrac{1}{2}|AP|\cdot d=9$，解得 $k=\dfrac{3}{2}$ 或 $k=-\dfrac{3}{2}$。

故 $l$ 的方程为 $y=\dfrac{3}{2}x-3$ 或 $y=-\dfrac{3}{2}x+6$。
\end{enumerate}

\item （证明与计算略，关键步骤：建立空间直角坐标系，利用向量法求解。）
\begin{enumerate}[label=(\arabic*)]
\item 证明略（利用线面平行的判定定理）。
\item $AD=3$。
\end{enumerate}

\item 
\begin{enumerate}[label=(\arabic*)]
\item 当 $b=0$ 时，$f(x)=\ln\dfrac{x}{2-x}+ax$，$f'(x)=\dfrac{2}{x(2-x)}+a$。

由 $f'(x)\geqslant 0$ 得 $a\geqslant -\dfrac{2}{x(2-x)}$。当 $x\in(0,2)$ 时，$x(2-x)\leqslant 1$，故 $-\dfrac{2}{x(2-x)}\leqslant -2$。因此 $a$ 的最小值为 $-2$。

\item 证明略（利用中心对称的定义，验证 $f(x)+f(2-x)=2a$ 为常数）。

\item $b\in\left[-\dfrac{2}{3},+\infty\right)$。
\end{enumerate}

\item 
\begin{enumerate}[label=(\arabic*)]
\item $(i,j)=(1,2),(1,6),(5,6)$。
\item 证明略（构造分组：利用等差数列的性质，将剩余项按一定规律分为 $m$ 组，每组构成等比数列）。
\item 证明略（利用组合计数与不等式放缩）。
\end{enumerate}

\end{enumerate}

\end{document}